\hypertarget{bayes}{}
\hypertarget{multivariate-empirical-bayes-and-partial-pooling}{%
\chapter{Multivariate empirical Bayes and partial
pooling}\label{multivariate-empirical-bayes-and-partial-pooling}}

\hypertarget{the-winners-curse}{%
\section{The winner\textquotesingle s curse}\label{the-winners-curse}}

When the largest estimate is selected from 25 noisy estimates, it is
likely to be unusually high partly because of noise. Partial pooling
shrinks estimates toward a cross-sectional centre, with the amount of
shrinkage determined by sampling uncertainty and the estimated
dispersion of true alphas.

\begin{equation}
\hat{\alpha } | \alpha  \sim N(\alpha , \hat{V}_{\mathrm{HAC}}), \alpha  \sim N(\mu 1, \tau ^{2}I)
\end{equation}

\begin{equation}
E(\alpha  | data) = \mu 1 + \tau ^{2}(\hat{V}_{\mathrm{HAC}} + \tau ^{2}I)^{-1}(\hat{\alpha } - \mu 1)
\end{equation}

\begin{equation}
Var(\alpha  | data) = \tau ^{2}I - \tau ^{4}(\hat{V}_{\mathrm{HAC}} + \tau ^{2}I)^{-1}
\end{equation}

With a diagonal V, each portfolio shrinks independently. With the full
V, a portfolio\textquotesingle s posterior update also reflects how its
alpha-estimation error co-moves with the others.

\hypertarget{deriving-the-posterior-by-completing-the-square}{%
\section{Deriving the posterior by completing the
square}\label{deriving-the-posterior-by-completing-the-square}}

Write a = \ensuremath{\hat{\alpha }} and let 1 denote the N-vector of ones. Ignoring constants
that do not depend on \ensuremath{\alpha }, Bayes\textquotesingle{} rule gives:

\begin{equation}
log p(\alpha |a) = -\frac{1}{2}(a-\alpha ){}^{\prime}V^{-1}(a-\alpha ) - \frac{1}{2}(\alpha -\mu 1){}^{\prime}(\tau ^{-2}I)(\alpha -\mu 1) + C.
\end{equation}

Collect the quadratic and linear terms in \ensuremath{\alpha }:

\begin{equation}
log p(\alpha |a) = -\frac{1}{2} \alpha {}^{\prime}P\alpha  + \alpha {}^{\prime}b + C{}^{\prime}, P = V^{-1}+\tau ^{-2}I, b=V^{-1}a+\tau ^{-2}\mu 1.
\end{equation}

Complete the square with C\textsubscript{post}=P\textsuperscript{\ensuremath{-}1} and
m\textsubscript{post}=C\textsubscript{post}b:

\begin{equation}
-\frac{1}{2}\alpha {}^{\prime}P\alpha +\alpha {}^{\prime}b = -\frac{1}{2}(\alpha -m_{\mathrm{post}}){}^{\prime}P(\alpha -m_{\mathrm{post}}) + \frac{1}{2}m_{\mathrm{post}}{}^{\prime}Pm_{\mathrm{post}}.
\end{equation}

Therefore \ensuremath{\alpha }\textbar a is multivariate normal. Matrix identities turn the
precision-form solution into the update form shown above:

\begin{equation}
m_{\mathrm{post}}=\mu 1+\tau ^{2}(V+\tau ^{2}I)^{-1}(a-\mu 1),
\end{equation}

\begin{equation}
C_{\mathrm{post}}=(V^{-1}+\tau ^{-2}I)^{-1}=\tau ^{2}I-\tau ^{4}(V+\tau ^{2}I)^{-1}.
\end{equation}

If V is diagonal with entry s\textsubscript{i}\textsuperscript{2}, the formula reduces to
the familiar scalar weighted average:

\begin{equation}
E(\alpha _{i}|a_{i}) = w_{i}a_{i} + (1-w_{i})\mu , w_{i}=\tau ^{2}/(\tau ^{2}+s_{i}^{2}).
\end{equation}

For example, with \ensuremath{\mu } = \ensuremath{-}26 bps/year, \ensuremath{\tau } = 94 bps/year, a\textsubscript{i}
= 250 bps/year and s\textsubscript{i} = 150 bps/year,

\begin{equation}
w_{i}=94^{2}/(94^{2}+150^{2})=0.282, \quad \text{posterior mean}= 0.282(250)+0.718(-26)=51.8\ \mathrm{bps/year}.
\end{equation}

The estimate shrinks sharply because its sampling uncertainty exceeds
the cross-sectional prior dispersion. In the actual multivariate
calculation there is no independent scalar weight for each portfolio:
off-diagonal elements of V couple all updates.

\hypertarget{estimating-ux3bc-and-ux3c4-by-marginal-likelihood}{%
\section{Estimating \ensuremath{\mu } and \ensuremath{\tau } by marginal
likelihood}\label{estimating-ux3bc-and-ux3c4-by-marginal-likelihood}}

The marginal model is \ensuremath{\hat{\alpha }} \textasciitilde{} N(\ensuremath{\mu }1, V + \ensuremath{\tau }\textsuperscript{2}I). Profile
marginal maximum likelihood searches over \ensuremath{\tau } \ensuremath{\geq } 0, solves for \ensuremath{\mu }
conditional on \ensuremath{\tau }, and selects the pair that maximises the marginal
likelihood.

Let S(\ensuremath{\tau })=V+\ensuremath{\tau }\textsuperscript{2}I. The marginal log-likelihood is

\begin{equation}
\ell (\mu ,\tau )=-\frac{1}{2}[N log(2\pi )+log|S(\tau )|+(a-\mu 1){}^{\prime}S(\tau )^{-1}(a-\mu 1)].
\end{equation}

Differentiating with respect to \ensuremath{\mu } gives

\begin{equation}
\partial \ell /\partial \mu  = 1{}^{\prime}S^{-1}(a-\mu 1).
\end{equation}

Setting this to zero yields the generalised least-squares profile
estimate:

\begin{equation}
\hat{\mu }(\tau ) = [1{}^{\prime}S(\tau )^{-1}a]/[1{}^{\prime}S(\tau )^{-1}1].
\end{equation}

Substitute \ensuremath{\hat{\mu }}(\ensuremath{\tau }) back into \ensuremath{\ell } and optimise the resulting one-dimensional
function over \ensuremath{\tau } \ensuremath{\geq } 0. Optimising log \ensuremath{\tau } enforces positivity, but the \ensuremath{\tau } = 0
boundary must be evaluated explicitly rather than excluded by the
transformation.

\begin{verbatim}
def profiled_mu(alpha_hat, V, tau2):
    S = V + tau2 * np.eye(len(alpha_hat))
    one = np.ones(len(alpha_hat))
    Sinv_one = np.linalg.solve(S, one)
    return (Sinv_one @ alpha_hat) / (Sinv_one @ one)

def neg_profile_loglik(log_tau, alpha_hat, V):
    tau2 = np.exp(2 * log_tau)
    S = V + tau2 * np.eye(len(alpha_hat))
    mu = profiled_mu(alpha_hat, V, tau2)
    d = alpha_hat - mu
    sign, logdet = np.linalg.slogdet(S)
    if sign <= 0:
        return np.inf
    return 0.5 * (logdet + d @ np.linalg.solve(S, d))
\end{verbatim}

\hypertarget{raw-tests-and-posterior-intervals-answer-different-questions}{%
\section{Raw tests and posterior intervals answer different
questions}\label{raw-tests-and-posterior-intervals-answer-different-questions}}

\begin{longtable}[]{@{}p{0.20\textwidth}p{0.34\textwidth}p{0.36\textwidth}@{}}
\toprule\noalign{}
Object & Question & Uses \\
\midrule\noalign{}
\endhead
\bottomrule\noalign{}
\endlastfoot
Raw alpha p/q & Is the unpooled regression intercept distinguishable
from zero? & Testing and discovery control \\
Posterior mean & What is the partially pooled estimate under the
hierarchical model? & Uncertainty-aware scoring \\
Posterior interval & What range is plausible under the fitted
hierarchical model? & Magnitude uncertainty \\
Bootstrap rank interval & Where might the portfolio rank under
historical resampling? & Ranking uncertainty \\
\end{longtable}

\textbf{Paper result.} The fitted prior centre is \ensuremath{-}26 annualised bps and
prior dispersion is 94 bps. Partial pooling changes at least one
portfolio by ten rank positions. Posterior intervals remain approximate
because \ensuremath{\tau } uncertainty is not fully integrated and V is treated as known
during the conditional posterior calculation.

